\chapitre[Chapter 1: sets and maps. Solving small linear systems by substitution.]{2}{Équations linéaires et matrices}{Linear equations and matrices}
\resume{Vecteurs de $\R^n$ et de $\C^n$, produit scalaire et inégalité de Cauchy, Schwarz ; matrices lues par leurs colonnes, produit, inverse, indépendance ; puis l'élimination de Gauss écrite avec des matrices, qui aboutit à la factorisation $PA=LU$.}
\summary{Vectors in R^n and C^n, dot product and the Cauchy-Schwarz inequality; matrices read through their columns, product, inverse, independence; then Gaussian elimination written with matrices, leading to the factorisation PA = LU.}
\begin{objectifs}
\item Calculer produits scalaires, normes, angles ; démontrer l'inégalité de Cauchy, Schwarz.
\item Lire $Ax=b$ de deux façons : intersection d'hyperplans (lignes) et combinaison des colonnes (colonnes).
\item Multiplier des matrices de quatre façons et connaître les règles du calcul matriciel, notamment $AB\neq BA$.
\item Caractériser l'inversibilité d'une matrice carrée et calculer un inverse.
\item Factoriser $A=LU$ ou $PA=LU$ en accumulant explicitement les matrices d'élimination, puis résoudre $Ax=b$.
\end{objectifs}
\begin{prerequis}
Chapitre 1 (applications, bijections). Résolution de petits systèmes linéaires par substitution. Nombres complexes : module et conjugué.
\end{prerequis}

\section{Vecteurs de $\R^n$}
\begin{definition}[L'espace $\R^n$][réf. 2.1.1]
$\R^n$ est l'ensemble des $n$-uplets $u=(u_1,\dots,u_n)$ de réels, appelés \textbf{vecteurs} (ou points) ; les $u_i$ sont ses \textbf{composantes}, et $\mathbf 0=(0,\dots,0)$ est l'origine. On les additionne et on les multiplie par un réel $k$ (un \textbf{scalaire}) composante par composante :
\[u+v=(u_1+v_1,\dots,u_n+v_n),\qquad ku=(ku_1,\dots,ku_n).\]
Une \textbf{combinaison linéaire} de $u_1,\dots,u_r$ est un vecteur $k_1u_1+\dots+k_ru_r$. Dans les calculs matriciels, on écrit les vecteurs en \textbf{colonnes}.
\end{definition}
\begin{intuition}
Dans $\R^2$, les multiples $ku$ d'un vecteur non nul décrivent une droite passant par l'origine. Les combinaisons $cu+dv$ de deux vecteurs non nuls décrivent cette même droite si $u$ et $v$ sont colinéaires, et \emph{tout le plan} sinon : n'importe quel point s'atteint en marchant $c$ fois le long de $u$ puis $d$ fois le long de $v$. Toute l'algèbre linéaire part de cette question : \emph{quels vecteurs peut-on atteindre ?}
\end{intuition}
\begin{theoreme}[Les huit règles][réf. Thm 2.1.4]\label{thm-huit}
Pour $u,v,w\in\R^n$ et $k,l\in\R$ :
\begin{itemize}
\item pour l'addition : (1) $(u+v)+w=u+(v+w)$ ; (2) $u+v=v+u$ ; (3) $u+\mathbf 0=u$ ; (4) $u+(-u)=\mathbf 0$ ;
\item pour les scalaires : (5) $k(u+v)=ku+kv$ ; (6) $(k+l)u=ku+lu$ ; (7) $(kl)u=k(lu)$ ; (8) $1u=u$.
\end{itemize}
\end{theoreme}
\begin{preuve}
Chaque règle se vérifie composante par composante, où elle devient une propriété connue des réels. Par exemple la $i$-ème composante de $k(u+v)$ est $k(u_i+v_i)=ku_i+kv_i$, qui est la $i$-ème composante de $ku+kv$.
\end{preuve}
\begin{essentiel}
Ces huit règles sont exactement les axiomes qu'abstrait le chapitre 3 : tout ensemble muni d'une addition et d'une multiplication par les scalaires qui les vérifie sera appelé \emph{espace vectoriel}, et tout ce qu'on démontre à partir d'elles seules y restera vrai.
\end{essentiel}

\section{Produit scalaire, norme et distance}
\begin{definition}[Produit scalaire, norme, distance][réf. 2.1.2]
Pour $u,v\in\R^n$ :
\begin{itemize}
\item le \textbf{produit scalaire} est $u\cdot v=u_1v_1+\dots+u_nv_n$ ; $u$ et $v$ sont \textbf{orthogonaux} si $u\cdot v=0$ ;
\item la \textbf{norme} est $\norm u=\sqrt{u\cdot u}$ et la \textbf{distance} est $d(u,v)=\norm{u-v}$ ;
\item $w$ est \textbf{unitaire} si $\norm w=1$ ; pour $w\neq\mathbf 0$, le vecteur $w/\norm w$ est unitaire et de même direction ;
\item l'\textbf{angle} de deux vecteurs non nuls est le $\theta\in[0,\pi]$ tel que $\cos\theta=\dfrac{u\cdot v}{\norm u\norm v}$ (ce quotient est dans $[-1,1]$ grâce au Théorème \voir{thm-cs}).
\end{itemize}
\end{definition}
\begin{proposition}[Propriétés du produit scalaire][réf. Thm 2.1.8]\label{prop-ps}
Pour $u,v,w\in\R^n$ et $k\in\R$ : $(u+v)\cdot w=u\cdot w+v\cdot w$ ; $(ku)\cdot v=k(u\cdot v)$ ; $u\cdot v=v\cdot u$ ; $u\cdot u\ge0$, avec égalité si et seulement si $u=\mathbf 0$.
\end{proposition}
\begin{preuve}
Les trois premières se vérifient sur la formule. Enfin $u\cdot u=\sum u_i^2$ est une somme de carrés : elle est positive, et nulle seulement si chaque $u_i$ est nul.
\end{preuve}
\begin{proposition}[Pythagore][réf. Prop. 2.1.11]
Si $u\perp v$, alors $\norm{u-v}^2=\norm u^2+\norm v^2$ (et de même pour $u+v$).
\end{proposition}
\begin{preuve}
$\norm{u-v}^2=(u-v)\cdot(u-v)=u\cdot u-2\,u\cdot v+v\cdot v=\norm u^2+\norm v^2$ puisque $u\cdot v=0$.
\end{preuve}
\begin{theoreme}[Inégalité de Cauchy, Schwarz][réf. Thm 2.1.12]\label{thm-cs}
Pour tous $u,v\in\R^n$ : $\abs{u\cdot v}\le\norm u\,\norm v$. L'égalité a lieu si et seulement si $u$ et $v$ sont colinéaires.
\end{theoreme}
\begin{preuve}[Démonstration par normalisation]
Si $u=\mathbf 0$ ou $v=\mathbf 0$, les deux membres sont nuls. Sinon posons $a=u/\norm u$ et $b=v/\norm v$, unitaires. Pour des réels, $xy\le\frac12(x^2+y^2)$ car $(x-y)^2\ge0$ ; donc
\[a\cdot b=\sum a_ib_i\le\tfrac12\sum(a_i^2+b_i^2)=\tfrac12(1+1)=1.\]
Le même calcul avec $-a$ donne $-a\cdot b\le1$, d'où $\abs{a\cdot b}\le1$, c'est-à-dire $\abs{u\cdot v}\le\norm u\norm v$ après multiplication par $\norm u\norm v$.
\emph{Cas d'égalité.} $a\cdot b=1$ impose $a_i=b_i$ pour tout $i$ (chaque inégalité $a_ib_i\le\frac12(a_i^2+b_i^2)$ doit être une égalité), donc $a=b$ ; de même $a\cdot b=-1$ impose $a=-b$. Dans les deux cas $v$ est un multiple de $u$. Réciproquement, si $v=ku$, alors $\abs{u\cdot v}=\abs k\norm u^2=\norm u\norm v$.
\end{preuve}
\begin{proposition}[Norme][réf. Prop. 2.1.14]\label{prop-norme}
$\norm u\ge0$, avec $\norm u=0\iff u=\mathbf 0$ ; $\norm{ku}=\abs k\norm u$ ; et l'\textbf{inégalité triangulaire} $\norm{u+v}\le\norm u+\norm v$, avec égalité si et seulement si l'un des vecteurs est un multiple positif ou nul de l'autre.
\end{proposition}
\begin{preuve}
Les deux premiers points découlent de \voir{prop-ps}. Pour le troisième :
\[\norm{u+v}^2=\norm u^2+2\,u\cdot v+\norm v^2\le\norm u^2+2\norm u\norm v+\norm v^2=(\norm u+\norm v)^2,\]
par Cauchy, Schwarz ; on prend la racine carrée. L'égalité exige $u\cdot v=\norm u\norm v$, c'est-à-dire $a\cdot b=1$ dans la preuve précédente : $u$ et $v$ colinéaires de même sens.
\end{preuve}

\section{Vecteurs de $\C^n$}
\begin{remarque}[Rappels][réf. Prop. 2.2.1]
Pour $z=a+ib$ : $\overline z=a-ib$, $\abs z=\sqrt{z\overline z}=\sqrt{a^2+b^2}$. On a $\overline{z+w}=\overline z+\overline w$, $\overline{zw}=\overline z\,\overline w$, $\overline{\overline z}=z$, $\abs{zw}=\abs z\abs w$ et $\abs{z+w}\le\abs z+\abs w$.
\end{remarque}
\begin{definition}[Produit scalaire complexe][réf. Déf. 2.2.2, 2.2.3]
Sur $\C^n$ (vecteurs à composantes complexes, additionnés et multipliés par des complexes composante par composante), on pose
\[z\cdot w=\sum_{i=1}^n z_i\overline{w_i},\qquad\norm z=\sqrt{z\cdot z}=\sqrt{\textstyle\sum\abs{z_i}^2}.\]
Orthogonalité et distance se définissent comme dans le cas réel.
\end{definition}
\begin{attention}[Ce qui change sur $\C$]
Sans le conjugué, $z\cdot z$ ne serait pas une longueur : pour $z=(1,i)$, $1^2+i^2=0$ alors que $z\neq\mathbf 0$. Avec le conjugué, $z\cdot z=\sum\abs{z_i}^2$ est un réel positif. Le prix à payer : le produit n'est plus symétrique ni linéaire des deux côtés.
\[w\cdot z=\overline{z\cdot w},\qquad (kz)\cdot w=k\,(z\cdot w),\qquad z\cdot(kw)=\overline k\,(z\cdot w).\]
On dit qu'il est \textbf{sesquilinéaire}. Enfin $\C^n$ est un espace vectoriel sur $\C$ \emph{et} sur $\R$ : le même ensemble porte deux structures, de dimensions $n$ et $2n$ (chapitre 3).
\end{attention}

\section{Matrices et produit matriciel}
\subsection{Le point de vue des colonnes}
\begin{exemple}[Différences et sommes][réf. 2.3]\label{ex-differences}
Soient $u=\begin{pmatrix}1\\-1\\0\end{pmatrix}$, $v=\begin{pmatrix}0\\1\\-1\end{pmatrix}$, $w=\begin{pmatrix}0\\0\\1\end{pmatrix}$, et $A$ la matrice de colonnes $u,v,w$. Pour $x=(x_1,x_2,x_3)$,
\[Ax=x_1u+x_2v+x_3w=\begin{pmatrix}x_1\\x_2-x_1\\x_3-x_2\end{pmatrix}.\]
$A$ calcule des \emph{différences}. Résoudre $Ax=b$, c'est chercher la combinaison des colonnes égale à $b$ : on trouve $x_1=b_1$, $x_2=b_1+b_2$, $x_3=b_1+b_2+b_3$, des \emph{sommes}. Autrement dit $x=Sb$ avec $S=\begin{pmatrix}1&0&0\\1&1&0\\1&1&1\end{pmatrix}$, et l'on vérifie $AS=SA=I_3$ : $S$ est l'inverse de $A$.
\end{exemple}
\begin{essentiel}[Deux lectures de $Ax=b$]
\begin{itemize}
\item \textbf{Lignes} : chaque équation définit une droite (dans $\R^2$), un plan (dans $\R^3$) ; une solution est un point commun à toutes.
\item \textbf{Colonnes} : $Ax=x_1a_1+\dots+x_na_n$ où $a_j$ est la $j$-ème colonne ; résoudre, c'est écrire $b$ comme combinaison linéaire des colonnes. Ce point de vue est le plus fécond : $Ax=b$ a une solution si et seulement si $b$ est atteignable à partir des colonnes.
\end{itemize}
\end{essentiel}

\subsection{Règles de calcul}
\begin{definition}[Matrices][réf. 2.3.1]
Une \textbf{matrice} $m\times n$ à coefficients dans $\K$ est un tableau $A=(a_{ij})$ de $m$ lignes et $n$ colonnes ; formellement, une application $\ens{1,\dots,m}\times\ens{1,\dots,n}\to\K$. On note $\Mat_{m\times n}(\K)$ leur ensemble et $\Mat_n(\K)$ celui des matrices \textbf{carrées} d'ordre $n$. La somme (même taille) et le produit par un scalaire se font coefficient par coefficient. La \textbf{matrice identité} $I_n$ a des $1$ sur la diagonale et des $0$ ailleurs.
\end{definition}
\begin{definition}[Produit matriciel]\label{def-produit}
Si $A$ est $m\times n$ et $B$ est $n\times p$ (le nombre de colonnes de $A$ égale le nombre de lignes de $B$), $AB$ est la matrice $m\times p$ de coefficients
\[(AB)_{ij}=\sum_{k=1}^n a_{ik}b_{kj}=(\text{ligne } i \text{ de } A)\cdot(\text{colonne } j \text{ de } B).\]
\end{definition}
\begin{proposition}[Quatre lectures du produit][réf. 2.3.1, Cor. 2.3.12]\label{prop-quatre}
Notons $a_1,\dots,a_n$ les colonnes de $A$, $b_1,\dots,b_p$ les colonnes de $B$, $\alpha_i$ les lignes de $A$ et $\beta_k$ les lignes de $B$.
\begin{enumerate}
\item \textbf{Coefficient par coefficient} : $(AB)_{ij}=\alpha_i\cdot b_j$.
\item \textbf{Par colonnes} : $AB=(Ab_1\ \cdots\ Ab_p)$ ; chaque colonne de $AB$ est une combinaison des colonnes de $A$.
\item \textbf{Par lignes} : la ligne $i$ de $AB$ est $\alpha_iB$ ; chaque ligne de $AB$ est une combinaison des lignes de $B$.
\item \textbf{Colonnes fois lignes} : $AB=a_1\beta_1+a_2\beta_2+\dots+a_n\beta_n$, somme de $n$ matrices $m\times p$ de la forme (colonne)(ligne).
\end{enumerate}
\end{proposition}
\begin{preuve}
Toutes expriment la même somme $\sum_k a_{ik}b_{kj}$. Pour (2), la $j$-ème colonne de $AB$ a pour $i$-ème coefficient $\sum_k a_{ik}b_{kj}$, qui est la $i$-ème composante de $\sum_k b_{kj}a_k=Ab_j$. Pour (4), la matrice $a_k\beta_k$ a pour coefficient $(i,j)$ le produit $a_{ik}b_{kj}$ ; en sommant sur $k$, on retrouve $(AB)_{ij}$.
\end{preuve}
\begin{proposition}[Règles][réf. Prop. 2.3.6, 2.3.7]\label{prop-regles}
Lorsque les tailles le permettent : $A(B+C)=AB+AC$, $(A+B)C=AC+BC$, $k(AB)=(kA)B=A(kB)$, $(AB)C=A(BC)$, $I_mA=AI_n=A$, et $A(au+bv)=aAu+bAv$. En revanche, \textbf{en général $AB\neq BA$}.
\end{proposition}
\begin{preuve}[Associativité]
$((AB)C)_{ij}=\sum_l\left(\sum_k a_{ik}b_{kl}\right)c_{lj}=\sum_{k,l}a_{ik}b_{kl}c_{lj}=\sum_k a_{ik}\left(\sum_l b_{kl}c_{lj}\right)=(A(BC))_{ij}$. Contre-exemple à la commutativité : avec $A=\begin{pmatrix}0&1\\0&0\end{pmatrix}$ et $B=\begin{pmatrix}0&0\\1&0\end{pmatrix}$, $AB=\begin{pmatrix}1&0\\0&0\end{pmatrix}\neq BA=\begin{pmatrix}0&0\\0&1\end{pmatrix}$.
\end{preuve}
\begin{remarque}[Produit par blocs][réf. Prop. 2.3.11]
Si l'on découpe $A$ et $B$ en blocs de façon compatible (les colonnes de $A$ découpées comme les lignes de $B$), on multiplie les blocs comme des nombres, en respectant l'ordre : $\begin{pmatrix}A_1&A_2\end{pmatrix}\begin{pmatrix}B_1\\B_2\end{pmatrix}=A_1B_1+A_2B_2$. La lecture (4) est le cas de blocs d'une colonne et d'une ligne.
\end{remarque}
\begin{proposition}[Matrices qui commutent avec tout][réf. Rem. 2.3.9, Question 2.3.10]
Les matrices $kI_n$ commutent avec toutes les matrices de $\Mat_n(\K)$, et ce sont les seules (exercice \voir{exo-centre}).
\end{proposition}

\section{Inverses}
\begin{definition}[Inverses à gauche et à droite][réf. 2.3.2]
Soit $A\in\Mat_{m\times n}(\K)$. Une matrice $G$ ($n\times m$) est un \textbf{inverse à gauche} si $GA=I_n$, une matrice $D$ ($n\times m$) un \textbf{inverse à droite} si $AD=I_m$.
\end{definition}
\begin{proposition}[réf. Prop. 2.3.13]\label{prop-inverse-unique}
Si $A\in\Mat_n(\K)$ a un inverse à gauche $G$ et un inverse à droite $D$, alors $G=D$. Cet inverse est alors unique, noté $A^{-1}$, et $A$ est dite \textbf{inversible}. L'ensemble des matrices inversibles d'ordre $n$ est noté $\GL_n(\K)$.
\end{proposition}
\begin{preuve}
$G=GI_n=G(AD)=(GA)D=I_nD=D$, grâce à l'associativité. Si $G'$ est un autre inverse, le même calcul avec $D$ montre $G'=D=G$.
\end{preuve}
\begin{proposition}[réf. Prop. 2.3.14]\label{prop-produit-inverse}
Si $A,B\in\GL_n(\K)$, alors $AB\in\GL_n(\K)$ et $(AB)^{-1}=B^{-1}A^{-1}$.
\end{proposition}
\begin{preuve}
$(AB)(B^{-1}A^{-1})=A(BB^{-1})A^{-1}=AA^{-1}=I_n$, et de même $(B^{-1}A^{-1})(AB)=I_n$. L'ordre s'inverse, comme pour des chaussettes et des chaussures : on retire en dernier ce qu'on a mis en premier.
\end{preuve}
\begin{attention}[réf. Question 2.3.15]
La somme de deux matrices inversibles n'est pas inversible en général : $A=I_n$ et $B=-I_n$ donnent $A+B=0$.
\end{attention}
\begin{theoreme}[Inversibilité et systèmes][réf. Prop. 2.3.16]\label{thm-inv-systeme}
$A\in\Mat_n(\K)$ est inversible si et seulement si, pour tout $b\in\K^n$, le système $Ax=b$ a une solution \textbf{unique}.
\end{theoreme}
\begin{preuve}
($\Rightarrow$) $x=A^{-1}b$ est solution ; si $Ax=b$, alors $x=A^{-1}Ax=A^{-1}b$.
($\Leftarrow$) Soit $d_j$ l'unique solution de $Ax=e_j$, où $e_j$ est la $j$-ème colonne de $I_n$, et $D$ la matrice de colonnes $d_j$ : $AD=I_n$ par la lecture (2). Reste à voir $DA=I_n$. On a $A(DA-I_n)=(AD)A-A=0$ : chaque colonne $c$ de $DA-I_n$ vérifie $Ac=\mathbf 0$. Comme $\mathbf 0$ est une solution de $Ax=\mathbf 0$ et que la solution est unique, $c=\mathbf 0$. Donc $DA=I_n$.
\end{preuve}
\begin{methode}[Calculer $A^{-1}$][réf. Rem. 2.3.17]
On résout simultanément les $n$ systèmes $Ax=e_j$ : on écrit le tableau $(A\mid I_n)$, on fait l'élimination jusqu'à obtenir $(I_n\mid A^{-1})$ (méthode de Gauss, Jordan). Si l'élimination fait apparaître un pivot nul impossible à remplacer, $A$ n'est pas inversible. Pour $2\times2$ : si $ad-bc\neq0$,
\[\begin{pmatrix}a&b\\c&d\end{pmatrix}^{-1}=\frac1{ad-bc}\begin{pmatrix}d&-b\\-c&a\end{pmatrix}.\]
\end{methode}

\section{Dépendance et indépendance linéaires}
\begin{definition}[Famille liée, libre][réf. Déf. 2.3.18]\label{def-libre}
Des vecteurs $u_1,\dots,u_r\in\K^n$ sont \textbf{liés} (linéairement dépendants) s'il existe des scalaires $x_1,\dots,x_r$ \textbf{non tous nuls} avec $x_1u_1+\dots+x_ru_r=\mathbf 0$. Sinon ils sont \textbf{libres} (linéairement indépendants) : la seule combinaison nulle est celle dont tous les coefficients sont nuls.
\end{definition}
\begin{remarque}
Les colonnes de $A$ sont libres si et seulement si $Ax=\mathbf 0$ n'a que la solution $x=\mathbf 0$. Deux vecteurs sont liés si et seulement si l'un est multiple de l'autre ; une famille contenant $\mathbf 0$ est liée.
\end{remarque}
\begin{proposition}[réf. Prop. 2.3.21]\label{prop-singuliere}
Pour $A\in\Mat_n(\K)$, si $Ax=\mathbf 0$ a une solution non nulle, alors $A$ n'est pas inversible (on dit \textbf{singulière}). Et si $Ax=b$ a deux solutions distinctes $x\neq y$, alors toutes les $x+t(y-x)$, $t\in\K$, sont solutions.
\end{proposition}
\begin{preuve}
Si $A$ était inversible, $Ax=\mathbf 0$ entraînerait $x=A^{-1}\mathbf 0=\mathbf 0$. Ensuite $A(x+t(y-x))=b+t(b-b)=b$.
\end{preuve}
La réciproque (\og colonnes libres $\Rightarrow$ inversible \fg) demande un argument de comptage : on l'obtiendra au Théorème \voir{thm-caract-inv} grâce à l'élimination.

\section{Transposition et permutations}
\begin{definition}[Transposée, matrices remarquables][réf. 2.3.4]
La \textbf{transposée} de $A\in\Mat_{m\times n}$ est $A\tr\in\Mat_{n\times m}$, $(A\tr)_{ij}=a_{ji}$ : les lignes de $A$ deviennent les colonnes de $A\tr$. Une matrice carrée est \textbf{triangulaire supérieure} si $a_{ij}=0$ pour $i>j$, \textbf{inférieure} si $a_{ij}=0$ pour $i<j$, \textbf{symétrique} si $A\tr=A$, \textbf{antisymétrique} si $A\tr=-A$.
\end{definition}
\begin{proposition}[réf. Prop. 2.3.24, Rem. 2.3.26, Prop. 2.3.30]\label{prop-transposee}
$(A+B)\tr=A\tr+B\tr$, $(kA)\tr=kA\tr$, $(A\tr)\tr=A$, $(AB)\tr=B\tr A\tr$, et si $A$ est inversible, $(A\tr)^{-1}=(A^{-1})\tr$. Dans $\R^n$, $x\cdot y=x\tr y$ et $(Ax)\cdot y=x\cdot(A\tr y)$ ; dans $\C^n$, $x\cdot y=x\tr\overline y$. Pour toute $A$, $A\tr A$ est symétrique.
\end{proposition}
\begin{preuve}
$((AB)\tr)_{ij}=(AB)_{ji}=\sum_k a_{jk}b_{ki}=\sum_k(B\tr)_{ik}(A\tr)_{kj}=(B\tr A\tr)_{ij}$. Ensuite $A\tr(A^{-1})\tr=(A^{-1}A)\tr=I\tr=I$, et de même de l'autre côté. Enfin $(Ax)\tr y=x\tr A\tr y$ et $(A\tr A)\tr=A\tr(A\tr)\tr=A\tr A$.
\end{preuve}
\begin{definition}[Matrices de permutation][réf. 2.3.4]
$P_{ij}$ est la matrice obtenue en échangeant les lignes $i$ et $j$ de $I_n$. Un produit de telles matrices est une \textbf{matrice de permutation} ; c'est une matrice qui a exactement un $1$ dans chaque ligne et chaque colonne, et des $0$ ailleurs. On note $\mathcal P_n$ leur ensemble.
\end{definition}
\begin{proposition}[réf. Prop. 2.3.33, corrigée]\label{prop-permutation}
\begin{enumerate}
\item $PA$ permute les \emph{lignes} de $A$, $AP$ permute ses \emph{colonnes}. En particulier $P_{ij}A$ échange les lignes $i$ et $j$ de $A$.
\item Toute $P\in\mathcal P_n$ est inversible et $P^{-1}=P\tr\in\mathcal P_n$.
\item $\Card\mathcal P_n=n!$.
\item $P^{n!}=I_n$ : la plus petite puissance $k\ge1$ avec $P^k=I_n$ divise $n!$. (Il est faux en général que $P^n=I_n$.)
\end{enumerate}
\end{proposition}
\begin{preuve}
(1) par les lectures (3) et (2) du produit. (2) Les colonnes de $P$ sont des vecteurs $e_k$ distincts, donc $(P\tr P)_{ij}=e_{k_i}\cdot e_{k_j}$ vaut $1$ si $i=j$ et $0$ sinon. (3) Une telle matrice est déterminée par la position du $1$ dans chaque ligne, positions deux à deux distinctes : c'est une bijection de $\ens{1,\dots,n}$, et il y en a $n!$ (\voir{prop-denombrement}). (4) $P$ permute les vecteurs $e_1,\dots,e_n$ selon une permutation dont les cycles ont des longueurs $\ell_1,\dots,\ell_r\le n$ (chapitre 4) ; $P^k=I_n$ dès que $k$ est un multiple commun des $\ell_i$, et $n!$ en est un. Contre-exemple : la matrice $4\times4$ qui permute circulairement $e_1,e_2,e_3$ et fixe $e_4$ vérifie $P^3=I_4$, donc $P^4=P\neq I_4$.
\end{preuve}

\section{Élimination et factorisation $LU$}
\begin{intuition}
La méthode du pivot de Gauss remplace une ligne par \og elle-même moins un multiple d'une ligne au-dessus \fg{} jusqu'à obtenir un système triangulaire. Chacune de ces opérations est une multiplication à gauche par une matrice très simple. En gardant la trace de ces matrices au lieu de les oublier, on obtient une factorisation de $A$ qui résume tout le calcul, et qu'on peut réutiliser pour chaque nouveau second membre.
\end{intuition}
\begin{definition}[Matrice d'élimination][réf. 2.4, Prop. 2.4.2]
Soient $k<i$ et $\ell\in\K$. La matrice $E_{ik}(\ell)$ est $I_n$ dans laquelle le coefficient $(i,k)$ est remplacé par $-\ell$. Multiplier $A$ à gauche par $E_{ik}(\ell)$ ne change que la ligne $i$, remplacée par $L_i-\ell L_k$. Si $a_{kj}\neq0$ est le pivot de la colonne $j$, le choix $\ell=a_{ij}/a_{kj}$ (le \textbf{multiplicateur}) met un $0$ en position $(i,j)$.
\end{definition}
\begin{proposition}[réf. Prop. 2.4.3]\label{prop-elim-inverse}
$E_{ik}(\ell)$ est inversible et $E_{ik}(\ell)^{-1}=E_{ik}(-\ell)$ : on change le signe de l'unique coefficient hors diagonale.
\end{proposition}
\begin{preuve}
$E_{ik}(-\ell)E_{ik}(\ell)A$ retranche $\ell L_k$ de la ligne $i$ puis l'ajoute ; la ligne $k$ n'a pas bougé entre-temps puisque $k\neq i$. Le produit est donc l'identité.
\end{preuve}
\begin{exemple}[Le fil rouge][réf. 2.4]\label{ex-fil-rouge}
Soit $A=\begin{pmatrix}2&1&1\\4&-6&0\\-2&7&2\end{pmatrix}$.
\begin{enumerate}
\item Pivot $2$. $\ell_{21}=4/2=2$ : $E_{21}A$ a pour deuxième ligne $(0,-8,-2)$.
\item $\ell_{31}=-2/2=-1$ : $E_{31}E_{21}A$ a pour troisième ligne $(-2,7,2)+(2,1,1)=(0,8,3)$.
\item Pivot $-8$. $\ell_{32}=8/(-8)=-1$ : la troisième ligne devient $(0,8,3)+(0,-8,-2)=(0,0,1)$.
\end{enumerate}
\[E_{32}E_{31}E_{21}A=U=\begin{pmatrix}2&1&1\\0&-8&-2\\0&0&1\end{pmatrix},\qquad A=E_{21}^{-1}E_{31}^{-1}E_{32}^{-1}U=LU,\quad L=\begin{pmatrix}1&0&0\\2&1&0\\-1&-1&1\end{pmatrix}.\]
Les multiplicateurs se rangent \emph{tels quels} sous la diagonale de $L$ : c'est le point clé.
\end{exemple}
\begin{proposition}[Pourquoi $L$ contient les multiplicateurs]
Si l'on élimine colonne par colonne, $L=E_{21}^{-1}E_{31}^{-1}\cdots E_{n,n-1}^{-1}$ est triangulaire inférieure à diagonale de $1$ et son coefficient $(i,k)$ est le multiplicateur $\ell_{ik}$.
\end{proposition}
\begin{preuve}
Lisons le produit de droite à gauche comme une suite d'opérations sur les lignes de $I_n$ : $E_{ik}^{-1}=E_{ik}(-\ell_{ik})$ \emph{ajoute} $\ell_{ik}$ fois la ligne $k$ à la ligne $i$. En partant de $I_n$ et en appliquant d'abord la dernière, $E_{n,n-1}^{-1}$, puis les précédentes, chaque opération ajoute un multiple d'une ligne $k$ qui, à ce moment, est encore la ligne $e_k\tr$ de l'identité (elle n'est modifiée que par des opérations d'indice de colonne plus petit, effectuées plus tard). Chaque multiplicateur se dépose donc en position $(i,k)$ sans interagir avec les autres. Dans l'ordre inverse, le produit $E_{32}E_{31}E_{21}$ mélange les multiplicateurs : son coefficient $(3,1)$ vaut $-\ell_{31}+\ell_{32}\ell_{21}$.
\end{preuve}
\begin{remarque}[Pivot nul][réf. Rem. 2.4.4]
Si le coefficient à éliminer est déjà nul, on passe au suivant. Si le \emph{pivot} est nul, on échange la ligne du pivot avec une ligne située en dessous qui a un coefficient non nul dans cette colonne (multiplication par une $P_{ij}$). On n'échange \textbf{jamais} de lignes quand le pivot est non nul. Si toute la colonne est nulle sous le pivot, il n'y a rien à éliminer : on passe à la colonne suivante, et $U$ aura un pivot nul.
\end{remarque}
\begin{theoreme}[Factorisation $PA=LU$][réf. Thm 2.4.5]\label{thm-lu}
Pour toute $A\in\Mat_n(\K)$, il existe une matrice de permutation $P$, une matrice $L$ triangulaire inférieure à diagonale de $1$ et une matrice $U$ triangulaire supérieure telles que $PA=LU$. Si aucun échange n'est nécessaire, $P=I_n$ et $A=LU$.
\end{theoreme}
\begin{preuve}[Idée de la démonstration]
On applique l'élimination en faisant tous les échanges de lignes \og à l'avance \fg : si on connaissait $P$, la matrice $PA$ s'éliminerait sans échange, et l'exemple précédent montre que cela donne $PA=LU$. Concrètement, quand un échange intervient au milieu du calcul, on échange aussi les multiplicateurs déjà rangés dans les lignes correspondantes de $L$ ; la figure ci-dessous fait exactement cela.
\end{preuve}
\widget{lu}
\begin{methode}[Résoudre $Ax=b$ avec $PA=LU$]
$Ax=b\iff PAx=Pb\iff L(Ux)=Pb$. On résout d'abord $Ly=Pb$ par \emph{descente} ($y_1$, puis $y_2$, ...), puis $Ux=y$ par \emph{remontée} ($x_n$, puis $x_{n-1}$, ...). La factorisation se calcule une fois (environ $n^3/3$ opérations) ; chaque nouveau second membre ne coûte ensuite qu'environ $n^2$ opérations.
\end{methode}
\begin{exemple}
Avec le fil rouge et $b=(5,-2,9)$ : $Ly=b$ donne $y_1=5$, $y_2=-2-2\cdot5=-12$, $y_3=9+5-12=2$. Puis $Ux=y$ : $x_3=2$, $-8x_2-4=-12$ donc $x_2=1$, $2x_1+1+2=5$ donc $x_1=1$. Contrôle : $A(1,1,2)\tr=(5,-2,9)\tr$.
\end{exemple}
\begin{theoreme}[Caractérisations de l'inversibilité][réf. Prop. 2.3.19, Thm 2.3.20]\label{thm-caract-inv}
Pour $A\in\Mat_n(\K)$, les propriétés suivantes sont équivalentes :
\begin{enumerate}
\item $A$ est inversible ;
\item pour tout $b$, $Ax=b$ a une solution unique ;
\item $Ax=\mathbf 0$ n'a que la solution $x=\mathbf 0$, c'est-à-dire les colonnes de $A$ sont libres ;
\item l'élimination produit $n$ pivots non nuls ;
\item les lignes de $A$ sont libres.
\end{enumerate}
\end{theoreme}
\begin{preuve}
(1)$\iff$(2) est le Théorème \voir{thm-inv-systeme}, et (2)$\Rightarrow$(3) est clair. (3)$\Rightarrow$(4) : écrivons $PA=LU$. Comme $P$ et $L$ sont inversibles, $Ax=\mathbf 0\iff Ux=\mathbf 0$. Si une colonne $j$ n'avait pas de pivot, on fixerait $x_j=1$, les autres variables sans pivot à $0$, et on calculerait les variables à pivot par remontée : on obtiendrait une solution non nulle de $Ux=\mathbf 0$. (4)$\Rightarrow$(2) : avec $n$ pivots non nuls, $Ly=Pb$ puis $Ux=y$ ont chacun une solution unique, obtenue par descente puis remontée. Enfin (1)$\iff$(5) : les lignes de $A$ sont les colonnes de $A\tr$, et $A$ est inversible si et seulement si $A\tr$ l'est (\voir{prop-transposee}).
\end{preuve}
\begin{remarque}[réf. Rem. 2.4.6]
Si $A$ est inversible, $A^{-1}=U^{-1}L^{-1}P$ : les inverses de matrices triangulaires se calculent par substitution, et $L^{-1}=E_{32}E_{31}E_{21}$ (dans cet ordre) est un produit de matrices d'élimination.
\end{remarque}

\begin{erreurs}
\item Multiplier dans le mauvais ordre ; supposer $AB=BA$ ; écrire $(AB)^{-1}=A^{-1}B^{-1}$ ou $(AB)\tr=A\tr B\tr$.
\item Permuter des lignes alors que le pivot n'est pas nul.
\item Oublier de changer le signe en inversant une matrice d'élimination.
\item Confondre inverse à gauche et inverse à droite pour des matrices non carrées.
\item Ranger dans $L$ les coefficients de $E_{32}E_{31}E_{21}$ au lieu des multiplicateurs.
\end{erreurs}

\exercices
\rubrique{A}
\begin{exo}[S][Produit scalaire, norme, angle]
Soient $u=(1,2,2)$, $v=(3,0,-4)$ et $w=(2,-1,0)$. Calculer $u\cdot v$, $\norm u$, $\norm v$, l'angle de $u$ et $v$, la distance $d(u,v)$, un vecteur unitaire de même direction que $u$. Les vecteurs $u$ et $w$ sont-ils orthogonaux ? Vérifier Pythagore pour $u$ et $w$.
\begin{corrige}
$u\cdot v=3+0-8=-5$, $\norm u=\sqrt{1+4+4}=3$, $\norm v=\sqrt{9+16}=5$, $\cos\theta=-\frac5{15}=-\frac13$, $\theta=\arccos(-\frac13)\approx109{,}47^\circ$ (angle obtus, produit négatif). $u-v=(-2,2,6)$, $d(u,v)=\sqrt{44}=2\sqrt{11}$. Unitaire : $\frac13(1,2,2)$. $u\cdot w=2-2+0=0$ : orthogonaux. $u-w=(-1,3,2)$, $\norm{u-w}^2=14=9+5=\norm u^2+\norm w^2$.
\end{corrige}
\end{exo}
\begin{exo}[S][Quatre façons de multiplier]
Soient $A=\begin{pmatrix}1&2&0\\0&1&-1\end{pmatrix}$ et $B=\begin{pmatrix}1&0\\2&1\\-1&3\end{pmatrix}$. Calculer $AB$ par les quatre lectures de la Proposition \voir{prop-quatre}. Le produit $BA$ existe-t-il ? Quelle est sa taille ?
\begin{corrige}
(1) $(AB)_{11}=1+4+0=5$, $(AB)_{12}=0+2+0=2$, $(AB)_{21}=0+2+1=3$, $(AB)_{22}=0+1-3=-2$ : $AB=\begin{pmatrix}5&2\\3&-2\end{pmatrix}$.
(2) Colonne 1 : $A(1,2,-1)\tr=1\cdot(1,0)\tr+2\cdot(2,1)\tr-1\cdot(0,-1)\tr=(5,3)\tr$ ; colonne 2 : $A(0,1,3)\tr=(2,1)\tr+3(0,-1)\tr=(2,-2)\tr$.
(3) Ligne 1 : $1\cdot(1,0)+2\cdot(2,1)+0\cdot(-1,3)=(5,2)$ ; ligne 2 : $0\cdot(1,0)+1\cdot(2,1)-1\cdot(-1,3)=(3,-2)$.
(4) $\begin{pmatrix}1\\0\end{pmatrix}(1\ 0)+\begin{pmatrix}2\\1\end{pmatrix}(2\ 1)+\begin{pmatrix}0\\-1\end{pmatrix}(-1\ 3)=\begin{pmatrix}1&0\\0&0\end{pmatrix}+\begin{pmatrix}4&2\\2&1\end{pmatrix}+\begin{pmatrix}0&0\\1&-3\end{pmatrix}=\begin{pmatrix}5&2\\3&-2\end{pmatrix}$.
$BA$ existe ($B$ a 2 colonnes, $A$ a 2 lignes) et est $3\times3$ : $AB$ et $BA$ n'ont même pas la même taille.
\end{corrige}
\end{exo}
\begin{exo}[S][Lignes et colonnes]
On considère le système $2x-y=1$, $x+y=5$. L'écrire sous la forme $Ax=b$, le résoudre, puis interpréter la solution (a) comme intersection de deux droites, (b) comme combinaison des colonnes de $A$. Faire un dessin pour chacune des deux lectures.
\begin{corrige}
$A=\begin{pmatrix}2&-1\\1&1\end{pmatrix}$, $b=(1,5)\tr$. En additionnant : $3x=6$, $x=2$, puis $y=3$. (a) Les droites $y=2x-1$ et $y=-x+5$ se coupent en $(2,3)$. (b) $2\begin{pmatrix}2\\1\end{pmatrix}+3\begin{pmatrix}-1\\1\end{pmatrix}=\begin{pmatrix}1\\5\end{pmatrix}$ : on atteint $b$ en marchant deux fois le long de la première colonne puis trois fois le long de la seconde.
\end{corrige}
\end{exo}
\begin{exo}[S][Combinaison linéaire ou non]
Soient $u_1=(1,0,1)$, $u_2=(0,1,2)$, $u_3=(1,1,3)$. (a) Le vecteur $b=(1,4,7)$ est-il combinaison linéaire de $u_1,u_2,u_3$ ? (b) Et $b'=(1,4,9)$ ? Donner toutes les écritures possibles. (c) Décrire géométriquement l'ensemble de toutes les combinaisons de $u_1,u_2,u_3$.
\begin{corrige}
On remarque $u_3=u_1+u_2$. (a) On résout $x_1+x_3=1$, $x_2+x_3=4$, $x_1+2x_2+3x_3=7$ : les deux premières donnent $x_1=1-x_3$, $x_2=4-x_3$, et la troisième devient $9=7$ : impossible. $b$ n'est pas combinaison. (b) La troisième équation devient $9=9$ : les solutions sont $(1-t,4-t,t)$, $t\in\R$ ; infinité d'écritures, car la famille est liée. (c) Les combinaisons de $u_1,u_2,u_3$ sont celles de $u_1,u_2$ (non colinéaires) : c'est le plan passant par $\mathbf 0$ engendré par $u_1$ et $u_2$, d'équation $z=x+2y$ (on vérifie $b'$ : $9=1+8$, et $b$ : $7\neq9$).
\end{corrige}
\end{exo}
\begin{exo}[S][Ce que l'on atteint]
Décrire géométriquement l'ensemble des combinaisons linéaires (a) de $u=(1,2,3)$ ; (b) de $u$ et $v=(2,4,6)$ ; (c) de $u$ et $w=(1,0,0)$ ; (d) de $u$, $w$ et $e_3=(0,0,1)$. Dans chaque cas, à quelle condition un vecteur $(x,y,z)$ est-il atteint ?
\begin{corrige}
(a) La droite $\R u$ : $(x,y,z)$ est atteint si et seulement si $y=2x$ et $z=3x$. (b) La même droite, car $v=2u$. (c) Le plan engendré par $u$ et $w$ ; un vecteur normal est $u\wedge w=(0,3,-2)$, d'où l'équation $3y-2z=0$. (d) Tout $\R^3$ : la matrice de colonnes $u,w,e_3$ a $3$ pivots non nuls (après échange, $w$, $u$, $e_3$ forment une matrice triangulaire de diagonale $1,2,1$), donc tout $b$ est atteint (\voir{thm-caract-inv}).
\end{corrige}
\end{exo}
\begin{exo}[S][Transposées et permutations]
Soit $A=\begin{pmatrix}1&2&3\\4&5&6\\7&8&9\end{pmatrix}$ et $P$ la matrice de permutation qui échange les lignes $1$ et $3$ de $I_3$. Calculer $A\tr$, $PA$, $AP$, $P\tr$, $P^2$. Écrire la matrice $Q$ qui envoie la ligne 1 en position 2, la ligne 2 en position 3 et la ligne 3 en position 1, et calculer $Q^3$.
\begin{corrige}
$A\tr=\begin{pmatrix}1&4&7\\2&5&8\\3&6&9\end{pmatrix}$. $P=\begin{pmatrix}0&0&1\\0&1&0\\1&0&0\end{pmatrix}$ ; $PA$ échange les lignes 1 et 3 : $\begin{pmatrix}7&8&9\\4&5&6\\1&2&3\end{pmatrix}$ ; $AP$ échange les colonnes 1 et 3 : $\begin{pmatrix}3&2&1\\6&5&4\\9&8&7\end{pmatrix}$. $P\tr=P$ et $P^2=I_3$. $Q=\begin{pmatrix}0&0&1\\1&0&0\\0&1&0\end{pmatrix}$ (la ligne 2 de $QA$ est la ligne 1 de $A$, etc.). Appliquer trois fois le décalage circulaire ramène chaque ligne à sa place : $Q^3=I_3$, alors que $Q^2=Q\tr\neq I_3$.
\end{corrige}
\end{exo}
\begin{exo}[S][Élimination et inverse d'une élimination]
Écrire la matrice $E$ qui, multipliée à gauche, retranche $3$ fois la ligne $1$ de la ligne $3$ d'une matrice $3\times3$. Écrire $E^{-1}$ et vérifier $EE^{-1}=I_3$. Calculer $E_{21}(2)E_{32}(5)$ et $E_{32}(5)E_{21}(2)$ ; lequel des deux produits contient \og les multiplicateurs à leur place \fg ?
\begin{corrige}
$E=\begin{pmatrix}1&0&0\\0&1&0\\-3&0&1\end{pmatrix}$, $E^{-1}=\begin{pmatrix}1&0&0\\0&1&0\\3&0&1\end{pmatrix}$ (changement de signe). Avec $E_{21}(2)$ (coefficient $-2$ en $(2,1)$) et $E_{32}(5)$ (coefficient $-5$ en $(3,2)$) : $E_{21}E_{32}=\begin{pmatrix}1&0&0\\-2&1&0\\0&-5&1\end{pmatrix}$ alors que $E_{32}E_{21}=\begin{pmatrix}1&0&0\\-2&1&0\\10&-5&1\end{pmatrix}$. Le premier range les coefficients à leur place ; c'est l'ordre \og $E_{21}^{-1}E_{31}^{-1}E_{32}^{-1}$ \fg{} qui sert pour $L$ (avec les signes changés).
\end{corrige}
\end{exo}
\begin{exo}[S][$A=LU$ et $PA=LU$]
(a) Factoriser $A=\begin{pmatrix}1&2&1\\2&3&1\\-1&1&4\end{pmatrix}$ sous la forme $A=LU$ en écrivant chaque matrice d'élimination. En déduire la solution de $Ax=(4,5,6)\tr$.
(b) Factoriser $B=\begin{pmatrix}0&2&1\\1&1&1\\2&0&3\end{pmatrix}$ sous la forme $PB=LU$. Pourquoi un échange est-il nécessaire ? Contrôler avec la figure (préréglage \og Pivot nul \fg).
\begin{corrige}
(a) $\ell_{21}=2$ : ligne 2 $\to(0,-1,-1)$. $\ell_{31}=-1$ : ligne 3 $\to(0,3,5)$. $\ell_{32}=3/(-1)=-3$ : ligne 3 $\to(0,3,5)+3(0,-1,-1)=(0,0,2)$. Donc $U=\begin{pmatrix}1&2&1\\0&-1&-1\\0&0&2\end{pmatrix}$, $L=\begin{pmatrix}1&0&0\\2&1&0\\-1&-3&1\end{pmatrix}$. $Ly=(4,5,6)$ : $y=(4,-3,6+4-9)=(4,-3,1)$. $Ux=y$ : $x_3=\frac12$, $-x_2-\frac12=-3$ donc $x_2=\frac52$, $x_1=4-5-\frac12=-\frac32$. Contrôle ligne 3 : $\frac32+\frac52+2=6$.
(b) Le pivot $b_{11}$ est nul : on échange les lignes 1 et 2 ($P=P_{12}$), ce qui donne $\begin{pmatrix}1&1&1\\0&2&1\\2&0&3\end{pmatrix}$. $\ell_{31}=2$ : ligne 3 $\to(0,-2,1)$ ; $\ell_{32}=-1$ : ligne 3 $\to(0,0,2)$. $U=\begin{pmatrix}1&1&1\\0&2&1\\0&0&2\end{pmatrix}$, $L=\begin{pmatrix}1&0&0\\0&1&0\\2&-1&1\end{pmatrix}$. Sans échange, aucune matrice d'élimination ne peut faire apparaître un pivot non nul en position $(1,1)$, car elles ne modifient que les lignes situées sous le pivot.
\end{corrige}
\end{exo}
\begin{exo}[S][Calculer des inverses]
(a) Inverser $\begin{pmatrix}2&1\\5&3\end{pmatrix}$.
(b) Inverser la matrice des différences $A=\begin{pmatrix}1&0&0\\-1&1&0\\0&-1&1\end{pmatrix}$ en résolvant $Ax=e_j$ pour $j=1,2,3$. Généraliser à $n\times n$.
(c) Inverser $\begin{pmatrix}1&a&b\\0&1&c\\0&0&1\end{pmatrix}$.
\begin{corrige}
(a) $ad-bc=1$ : l'inverse est $\begin{pmatrix}3&-1\\-5&2\end{pmatrix}$.
(b) $Ax=e_1$ : $x_1=1$, $x_2=x_1=1$, $x_3=x_2=1$ ; $Ax=e_2$ : $x=(0,1,1)$ ; $Ax=e_3$ : $x=(0,0,1)$. Donc $A^{-1}=S=\begin{pmatrix}1&0&0\\1&1&0\\1&1&1\end{pmatrix}$, l'exemple \voir{ex-differences}. En taille $n$, l'inverse de la matrice des différences est la matrice triangulaire inférieure remplie de $1$ : résoudre $Ax=b$ revient à faire des sommes cumulées.
(c) On résout par remontée, ou l'on vérifie : $\begin{pmatrix}1&-a&ac-b\\0&1&-c\\0&0&1\end{pmatrix}$. Le coefficient $(1,3)$ vaut $ac-b$ et non $-b$ : l'inverse d'un produit d'éliminations n'est pas \og le même avec les signes changés \fg{} en dehors de l'ordre favorable.
\end{corrige}
\end{exo}
\begin{exo}[A][Pour quelles valeurs ?]
Pour quelles valeurs de $a$ la matrice $\begin{pmatrix}a&1&1\\1&a&1\\1&1&a\end{pmatrix}$ est-elle singulière ? Raisonner par élimination, en évitant de diviser par une quantité qui pourrait être nulle.
\begin{indication}
Échange d'abord les lignes 1 et 3, pour avoir un pivot égal à $1$ qui ne dépend pas de $a$.
\end{indication}
\begin{corrige}
Avec les lignes dans l'ordre 3, 2, 1 : $\begin{pmatrix}1&1&a\\1&a&1\\a&1&1\end{pmatrix}$. $L_2-L_1=(0,a-1,1-a)$, $L_3-aL_1=(0,1-a,1-a^2)$. Si $a=1$, ces deux lignes sont nulles : la matrice est singulière. Si $a\neq1$, le pivot $a-1$ est non nul et $L_3+L_2=(0,0,(1-a)+(1-a^2))=(0,0,(1-a)(2+a))$. Les pivots sont $1$, $a-1$, $(1-a)(a+2)$. La matrice est singulière si et seulement si $a=1$ ou $a=-2$. (Contrôle : pour $a=-2$, la somme des trois lignes est nulle, et pour $a=1$, les trois lignes sont égales.)
\end{corrige}
\end{exo}
\begin{exo}[A][Symétrique plus antisymétrique]
Montrer que toute $A\in\Mat_n(\R)$ s'écrit de manière unique $A=S+T$ avec $S$ symétrique et $T$ antisymétrique, et exprimer $S$ et $T$ à l'aide de $A$ et $A\tr$. Application : $A=\begin{pmatrix}1&2\\4&3\end{pmatrix}$.
\begin{corrige}
\emph{Analyse.} Si $A=S+T$, alors $A\tr=S-T$, donc nécessairement $S=\frac12(A+A\tr)$ et $T=\frac12(A-A\tr)$ : d'où l'unicité. \emph{Synthèse.} Ces deux matrices conviennent : $S\tr=\frac12(A\tr+A)=S$, $T\tr=\frac12(A\tr-A)=-T$ et $S+T=A$. Pour l'exemple : $S=\begin{pmatrix}1&3\\3&3\end{pmatrix}$, $T=\begin{pmatrix}0&-1\\1&0\end{pmatrix}$.
\end{corrige}
\end{exo}

\rubrique{B}
\begin{exo}[S][Cauchy, Schwarz et l'inégalité triangulaire]
(a) Rédiger la preuve de l'inégalité de Cauchy, Schwarz et de son cas d'égalité.
(b) Autre preuve : pour $u,v$ fixés avec $v\neq\mathbf 0$, étudier le signe du polynôme $t\mapsto\norm{u+tv}^2$.
(c) En déduire l'inégalité triangulaire et $\abs{\norm u-\norm v}\le\norm{u-v}$.
\begin{corrige}
(a) Voir \voir{thm-cs}.
(b) $\norm{u+tv}^2=\norm v^2t^2+2(u\cdot v)t+\norm u^2\ge0$ pour tout $t$ : ce trinôme de degré $2$ ne change pas de signe, donc son discriminant est négatif ou nul : $4(u\cdot v)^2-4\norm u^2\norm v^2\le0$. Égalité si et seulement si le trinôme a une racine $t_0$, c'est-à-dire $u+t_0v=\mathbf 0$.
(c) Voir \voir{prop-norme} pour la première. Pour la seconde, $\norm u=\norm{(u-v)+v}\le\norm{u-v}+\norm v$, donc $\norm u-\norm v\le\norm{u-v}$ ; en échangeant les rôles, $\norm v-\norm u\le\norm{u-v}$.
\end{corrige}
\end{exo}
\begin{exo}[S][Règles de calcul]
Démontrer, coefficient par coefficient ou par un calcul matriciel : (a) $(AB)\tr=B\tr A\tr$ ; (b) si $A,B$ sont inversibles, $(AB)^{-1}=B^{-1}A^{-1}$ ; (c) $(A\tr)^{-1}=(A^{-1})\tr$ ; (d) si $A$ est symétrique inversible, $A^{-1}$ est symétrique.
\begin{corrige}
(a), (b), (c) : voir \voir{prop-transposee} et \voir{prop-produit-inverse}. (d) $(A^{-1})\tr=(A\tr)^{-1}=A^{-1}$ par (c) et $A\tr=A$.
\end{corrige}
\end{exo}
\begin{exo}[S][L'inverse est unique]
Soit $A\in\Mat_n(\K)$. (a) Montrer que si $GA=I_n$ et $AD=I_n$, alors $G=D$. (b) Montrer que si $AD=I_n$, alors $A$ est inversible et $D=A^{-1}$. (On pourra utiliser le Théorème \voir{thm-caract-inv}.) (c) Trouver une matrice $2\times3$ qui a une infinité d'inverses à droite et aucun inverse à gauche.
\begin{corrige}
(a) $G=G(AD)=(GA)D=D$. (b) Si $Dx=\mathbf 0$, alors $x=ADx=\mathbf 0$ : les colonnes de $D$ sont libres, donc $D$ est inversible (\voir{thm-caract-inv}), et $A=AI_n=ADD^{-1}=D^{-1}$, inversible d'inverse $D$. Pour des matrices carrées, un inverse d'un seul côté suffit donc. (c) $A=\begin{pmatrix}1&0&0\\0&1&0\end{pmatrix}$ : $AD=I_2$ pour toute $D=\begin{pmatrix}1&0\\0&1\\a&b\end{pmatrix}$. Si $GA=I_3$, alors $GA(0,0,1)\tr=(0,0,1)\tr$ alors que $A(0,0,1)\tr=\mathbf 0$ : impossible.
\end{corrige}
\end{exo}
\begin{exo}[S][Un groupe de matrices]
Un \emph{groupe de matrices} est une partie $\mathcal G$ de $\GL_n(\K)$ stable par produit et par passage à l'inverse. (a) Montrer qu'un groupe de matrices non vide contient $I_n$. (b) Montrer que l'ensemble $\mathcal U$ des matrices $\begin{pmatrix}1&a&b\\0&1&c\\0&0&1\end{pmatrix}$ ($a,b,c\in\R$) est un groupe de matrices. Est-il commutatif ? (c) Les matrices de permutation $n\times n$ forment-elles un groupe de matrices ?
\begin{corrige}
(a) Si $M\in\mathcal G$, alors $M^{-1}\in\mathcal G$, puis $MM^{-1}=I_n\in\mathcal G$.
(b) Produit : $\begin{pmatrix}1&a&b\\0&1&c\\0&0&1\end{pmatrix}\begin{pmatrix}1&a'&b'\\0&1&c'\\0&0&1\end{pmatrix}=\begin{pmatrix}1&a+a'&b'+ac'+b\\0&1&c+c'\\0&0&1\end{pmatrix}\in\mathcal U$. Inverse : l'exercice précédent sur les inverses donne $\begin{pmatrix}1&-a&ac-b\\0&1&-c\\0&0&1\end{pmatrix}\in\mathcal U$ (et ces matrices sont bien inversibles). Non commutatif : le coefficient $(1,3)$ du produit vaut $b+b'+ac'$ dans un ordre et $b+b'+a'c$ dans l'autre ; avec $a=1,c'=1,a'=c=0$, on obtient $1\neq0$.
(c) Oui : $\mathcal P_n\subset\GL_n$ est stable par produit (par définition, produits de $P_{ij}$) et par inverse ($P^{-1}=P\tr\in\mathcal P_n$, \voir{prop-permutation}).
\end{corrige}
\end{exo}
\begin{exo}[S][Échanger $AB$ et $BA$]
Soient $A\in\Mat_{m\times n}$ et $B\in\Mat_{n\times m}$. (a) Vérifier $B(I_m-AB)=(I_n-BA)B$. (b) On suppose $I_m-AB$ inversible, d'inverse $C$. En cherchant l'inverse de $I_n-BA$ sous la forme $I_n+X$, \emph{trouver} par le calcul un candidat, puis vérifier qu'il est inverse des deux côtés. (c) En déduire que $I_m-AB$ est inversible si et seulement si $I_n-BA$ l'est.
\begin{indication}
L'équation $(I_n-BA)(I_n+X)=I_n$ se réécrit $(I_n-BA)X=BA$. Cherche $X$ sous la forme $BYA$ et utilise (a).
\end{indication}
\begin{corrige}
(a) Les deux membres valent $B-BAB$.
(b) $(I_n-BA)(I_n+X)=I_n\iff X-BAX=BA\iff(I_n-BA)X=BA$. Avec $X=BYA$ et (a) : $(I_n-BA)BYA=B(I_m-AB)YA$, qui vaut $BA$ dès que $(I_m-AB)Y=I_m$, c'est-à-dire $Y=C$. Candidat : $I_n+BCA$. Vérification à droite : fait par construction. À gauche : $(I_n+BCA)(I_n-BA)=I_n-BA+BCA-BCABA=I_n-BA+BC(I_m-AB)A=I_n-BA+BA=I_n$.
(c) Le sens direct est (b) ; la réciproque s'obtient en échangeant les rôles de $A$ et $B$.
\end{corrige}
\end{exo}
\begin{exo}[A][Triangulaire unitaire sans déterminant]
Soit $L$ une matrice triangulaire inférieure dont la diagonale ne contient que des $1$. Montrer, sans déterminant, que $L$ est inversible et que $L^{-1}$ est encore triangulaire inférieure à diagonale de $1$.
\begin{corrige}
Si $Lx=\mathbf 0$, la première équation donne $x_1=0$, puis la deuxième $\ell_{21}x_1+x_2=0$ donne $x_2=0$, et ainsi de suite (descente) : $x=\mathbf 0$. Les colonnes sont libres, donc $L$ est inversible (\voir{thm-caract-inv}). Pour la forme de $L^{-1}$ : l'élimination de $L$ n'utilise que des $E_{ik}$ avec $i>k$ et aucun échange, et aboutit à $U=I_n$ ; donc $L^{-1}$ est un produit de matrices d'élimination, chacune triangulaire inférieure à diagonale de $1$, et ce type de matrice est stable par produit (le coefficient $(i,j)$ d'un produit, $i<j$, est une somme de termes contenant un facteur nul).
\end{corrige}
\end{exo}
\begin{exo}[A][Le centre de $\Mat_n(\K)$]\label{exo-centre}
Soit $E_{ij}$ la matrice dont le seul coefficient non nul est un $1$ en position $(i,j)$. (a) Décrire $AE_{ij}$ et $E_{ij}A$. (b) Montrer que si $A$ commute avec toutes les matrices, alors $A=kI_n$.
\begin{corrige}
(a) $AE_{ij}$ a pour $j$-ème colonne la $i$-ème colonne de $A$ et des zéros ailleurs ; $E_{ij}A$ a pour $i$-ème ligne la $j$-ème ligne de $A$ et des zéros ailleurs.
(b) Si $AE_{ij}=E_{ij}A$, comparons. Le coefficient $(k,j)$ vaut $a_{ki}$ à gauche, et à droite il vaut $0$ si $k\neq i$ et $a_{jj}$ si $k=i$. Donc $a_{ki}=0$ pour $k\neq i$ (les coefficients hors diagonale sont nuls) et $a_{ii}=a_{jj}$ pour tous $i,j$. Ainsi $A=a_{11}I_n$.
\end{corrige}
\end{exo}
\begin{exo}[A][Ce qui change dans $\C^n$]
Reprendre les résultats de la section 2.2 dans $\C^n$ avec $z\cdot w=\sum z_i\overline{w_i}$ et dire précisément ce qui reste vrai : symétrie, linéarité, positivité, Pythagore, Cauchy, Schwarz, inégalité triangulaire.
\begin{corrige}
Positivité : $z\cdot z=\sum\abs{z_i}^2\ge0$, nul seulement pour $z=\mathbf 0$ : vrai. Symétrie : remplacée par $w\cdot z=\overline{z\cdot w}$. Linéarité à gauche vraie, à droite \og conjuguée \fg : $z\cdot(kw)=\overline k(z\cdot w)$. Pythagore : $\norm{z+w}^2=\norm z^2+2\,\mathrm{Re}(z\cdot w)+\norm w^2$, donc si $z\perp w$ l'égalité reste vraie ; mais la réciproque devient fausse : $\mathrm{Re}(z\cdot w)=0$ suffit (exemple $z=(1)$, $w=(i)$ dans $\C^1$). Cauchy, Schwarz reste vrai : pour $w\neq\mathbf 0$, poser $\lambda=\frac{z\cdot w}{\norm w^2}$ et développer $0\le\norm{z-\lambda w}^2=\norm z^2-\frac{\abs{z\cdot w}^2}{\norm w^2}$. L'inégalité triangulaire en découle comme dans le cas réel, puisque $\mathrm{Re}(z\cdot w)\le\abs{z\cdot w}$.
\end{corrige}
\end{exo}
\begin{exo}[A][Matrices de permutation]
Soit $P\in\mathcal P_n$. (a) Montrer que $P\tr P=I_n$. (b) Montrer que l'ensemble $\ens{P^k\mid k\ge1}$ est fini et en déduire qu'il existe $k\ge1$ avec $P^k=I_n$. (c) Trouver dans $\mathcal P_4$ une matrice d'ordre $4$ et une matrice d'ordre $3$. Peut-on trouver une matrice de $\mathcal P_4$ d'ordre $6$ ? et dans $\mathcal P_5$ ?
\begin{corrige}
(a) Voir \voir{prop-permutation}. (b) Les $P^k$ sont dans $\mathcal P_n$, qui a $n!$ éléments ; deux d'entre elles coïncident : $P^a=P^b$ avec $a<b$, et en multipliant par $P^{-a}$, $P^{b-a}=I_n$. (c) La matrice de la permutation circulaire $e_1\to e_2\to e_3\to e_4\to e_1$ est d'ordre $4$ ; celle de $e_1\to e_2\to e_3\to e_1$ (et $e_4$ fixe) d'ordre $3$. L'ordre est le ppcm des longueurs des cycles (chapitre 4) ; dans $\mathcal P_4$, les types de cycles possibles donnent les ordres $1,2,3,4$ : pas d'ordre $6$. Dans $\mathcal P_5$, un $3$-cycle et une transposition disjoints donnent l'ordre $6$.
\end{corrige}
\end{exo}

\rubrique{C}
\begin{exo}[S][Vrai ou faux sur les produits]
Pour des matrices carrées : (a) $AB=BA$ ; (b) $AB=0\Rightarrow A=0$ ou $B=0$ ; (c) $A^2=0\Rightarrow A=0$ ; (d) $A,B$ inversibles $\Rightarrow A+B$ inversible ; (e) $A,B$ inversibles $\Rightarrow AB$ inversible ; (f) $(A+B)^2=A^2+2AB+B^2$.
\begin{corrige}
(a) Faux en général (\voir{prop-regles}). (b) Faux : $A=\begin{pmatrix}1&0\\0&0\end{pmatrix}$, $B=\begin{pmatrix}0&0\\0&1\end{pmatrix}$. (c) Faux : $A=\begin{pmatrix}0&1\\0&0\end{pmatrix}$. (d) Faux : $B=-A$. (e) Vrai (\voir{prop-produit-inverse}). (f) Faux en général : $(A+B)^2=A^2+AB+BA+B^2$, et l'égalité équivaut à $AB=BA$.
\end{corrige}
\end{exo}
\begin{exo}[S][Combien de solutions ?]
(a) Un système linéaire peut-il avoir exactement deux solutions ? (b) Si $x$ et $y$ sont solutions de $Ax=b$, que dire de $x-y$, de $\frac{x+y}2$, de $x+y$ ? (c) Si $x_0$ est une solution de $Ax=b$, décrire toutes les solutions à l'aide des solutions de $Ax=\mathbf 0$.
\begin{corrige}
(a) Non : $0$, $1$ ou une infinité (sur $\R$ ou $\C$), car deux solutions distinctes en donnent une droite entière (\voir{prop-singuliere}). (b) $A(x-y)=\mathbf 0$ : solution du système homogène. $A\frac{x+y}2=b$ : solution. $A(x+y)=2b$ : solution seulement si $b=\mathbf 0$. (c) $Ax=b\iff A(x-x_0)=\mathbf 0$ : les solutions sont les $x_0+h$, $h$ parcourant les solutions de l'équation homogène. \og Solution générale = solution particulière + solution générale homogène \fg.
\end{corrige}
\end{exo}
\begin{exo}[S][Qui est stable ?]
Parmi les ensembles suivants de matrices $n\times n$ ($n\ge2$), lesquels sont stables par produit ? Parmi les inversibles qu'ils contiennent, lesquels sont stables par inverse ? Lesquels sont des groupes de matrices ? (a) triangulaires supérieures inversibles ; (b) symétriques inversibles ; (c) matrices de permutation ; (d) matrices orthogonales ($A\tr A=I_n$) ; (e) matrices à coefficients entiers inversibles.
\begin{corrige}
(a) Stable par produit (coefficient $(i,j)$, $i>j$, somme de termes contenant un zéro) et par inverse (remontée) : groupe. (b) Inverse symétrique, mais produit non symétrique en général : $\begin{pmatrix}1&0\\0&2\end{pmatrix}\begin{pmatrix}0&1\\1&0\end{pmatrix}=\begin{pmatrix}0&1\\2&0\end{pmatrix}$ : pas un groupe. (c) Groupe. (d) $(AB)\tr AB=B\tr A\tr AB=I_n$ et $A^{-1}=A\tr$ vérifie $(A\tr)\tr A\tr=AA\tr=I_n$ : groupe. (e) Stable par produit mais pas par inverse : $\begin{pmatrix}2&0\\0&1\end{pmatrix}^{-1}$ n'est pas à coefficients entiers. Les matrices entières dont l'inverse est entier forment le groupe $\GL_n(\Z)$ (chapitre 4).
\end{corrige}
\end{exo}
\begin{exo}[A][Angles obtus]
(a) Trouver trois vecteurs de $\R^2$ dont les produits scalaires deux à deux sont strictement négatifs. (b) Montrer qu'il n'en existe pas quatre. (c) Donner quatre vecteurs de $\R^3$ avec la même propriété.
\begin{corrige}
(a) Trois vecteurs unitaires à $120^\circ$ : $(1,0)$, $(-\frac12,\frac{\sqrt3}2)$, $(-\frac12,-\frac{\sqrt3}2)$ ; produits $-\frac12$. (b) Des vecteurs non nuls deux à deux à angle strictement obtus sont deux à deux distincts en direction. Rangeons quatre tels vecteurs selon leur angle polaire ; les quatre écarts entre angles consécutifs (en faisant le tour) totalisent $360^\circ$, et chacun doit dépasser $90^\circ$ : impossible. (c) Les sommets d'un tétraèdre régulier centré en $\mathbf 0$ : $(1,1,1)$, $(1,-1,-1)$, $(-1,1,-1)$, $(-1,-1,1)$, de produits deux à deux égaux à $-1$. En général on peut en trouver $n+1$ dans $\R^n$, et pas davantage.
\end{corrige}
\end{exo}
\begin{exo}[A][Le produit scalaire complexe]
Vrai ou faux dans $\C^n$ ? (a) $z\cdot w=w\cdot z$ ; (b) $(kz)\cdot w=z\cdot(kw)$ ; (c) $z\cdot z=0\Rightarrow z=\mathbf 0$ ; (d) $\sum z_iw_i=0$ pour $w=z$ entraîne $z=\mathbf 0$.
\begin{corrige}
(a) Faux : $w\cdot z=\overline{z\cdot w}$ ; avec $z=(i)$, $w=(1)$ : $z\cdot w=i$ et $w\cdot z=-i$. (b) Faux : à gauche $k(z\cdot w)$, à droite $\overline k(z\cdot w)$ ; égalité si $k$ est réel ou $z\cdot w=0$. (c) Vrai, grâce au conjugué. (d) Faux : $z=(1,i)$ donne $1+i^2=0$. C'est pourquoi on conjugue.
\end{corrige}
\end{exo}

\rubrique{D}
\begin{exo}[S][Une matrice de bout en bout]
Pour $A=\begin{pmatrix}2&1&1\\4&-6&0\\-2&7&2\end{pmatrix}$ : (a) retrouver $A=LU$ ; (b) résoudre $Ax=(5,-2,9)\tr$ par descente et remontée ; (c) calculer $A^{-1}$ par la méthode de Gauss, Jordan ; (d) vérifier que les quatre résultats sont cohérents : $A^{-1}b$ redonne la solution de (b), et le produit des pivots est non nul.
\begin{corrige}
(a) et (b) : voir l'exemple \voir{ex-fil-rouge} et la suite ; $x=(1,1,2)$.
(c) À partir de $(A\mid I_3)$ : les éliminations donnent $(U\mid L^{-1})$ avec $L^{-1}=\begin{pmatrix}1&0&0\\-2&1&0\\-1&1&1\end{pmatrix}$ (lignes : $L_2-2L_1$, puis $L_3+L_1$, puis $L_3+L_2$). La remontée (diviser $L_3$ par $1$, $L_2\leftarrow(L_2+2L_3)/(-8)$, $L_1\leftarrow(L_1-L_2-L_3)/2$) donne
\[A^{-1}=\begin{pmatrix}\frac34&-\frac5{16}&-\frac38\\[2pt]\frac12&-\frac38&-\frac14\\[2pt]-1&1&1\end{pmatrix}.\]
(d) $A^{-1}(5,-2,9)\tr$ : première ligne $\frac{15}4+\frac{10}{16}-\frac{27}8=\frac{60+10-54}{16}=1$ ; deuxième $\frac52+\frac34-\frac94=1$ ; troisième $-5-2+9=2$. On retrouve $(1,1,2)$. Les pivots $2,-8,1$ sont non nuls, ce qui garantissait l'inversibilité (\voir{thm-caract-inv}) ; leur produit $-16$ sera le déterminant (chapitre 3).
\end{corrige}
\end{exo}
\begin{exo}[A][$LU$ avec des paramètres]
Factoriser $A=\begin{pmatrix}a&a&a\\a&b&b\\a&b&c\end{pmatrix}$ sous la forme $LU$. À quelles conditions sur $a,b,c$ la matrice est-elle inversible ? Reconnaître $L$.
\begin{corrige}
Si $a\neq0$ : $\ell_{21}=\ell_{31}=1$, les lignes 2 et 3 deviennent $(0,b-a,b-a)$ et $(0,b-a,c-a)$. Si $b\neq a$ : $\ell_{32}=1$, la ligne 3 devient $(0,0,c-b)$. Donc $L=\begin{pmatrix}1&0&0\\1&1&0\\1&1&1\end{pmatrix}$ (la matrice des sommes $S$) et $U=\begin{pmatrix}a&a&a\\0&b-a&b-a\\0&0&c-b\end{pmatrix}$. $A$ est inversible si et seulement si les trois pivots $a$, $b-a$, $c-b$ sont non nuls. (Si $a=0$ ou $b=a$, la matrice a deux lignes égales ou une ligne nulle ; si $c=b$, deux lignes égales : elle est singulière, ce qui est cohérent.)
\end{corrige}
\end{exo}
\begin{exo}[A][Un problème de modélisation]
On cherche un polynôme $p(x)=a+bx+cx^2$ dont le graphe passe par $(1,2)$, $(2,3)$ et $(3,6)$. (a) Écrire le problème sous la forme $Ax=b$. (b) Montrer qu'il a une solution unique et la calculer. (c) Interpréter : en termes de plans de $\R^3$, puis en termes de combinaison des colonnes. (d) Montrer plus généralement que trois points d'abscisses distinctes déterminent un unique tel polynôme.
\begin{corrige}
(a) $\begin{pmatrix}1&1&1\\1&2&4\\1&3&9\end{pmatrix}\begin{pmatrix}a\\b\\c\end{pmatrix}=\begin{pmatrix}2\\3\\6\end{pmatrix}$.
(b) Élimination : $L_2-L_1=(0,1,3\mid1)$, $L_3-L_1=(0,2,8\mid4)$, puis $L_3-2L_2=(0,0,2\mid2)$. Trois pivots non nuls : solution unique. $c=1$, $b=1-3=-2$, $a=2+2-1=3$ : $p(x)=x^2-2x+3$.
(c) Trois plans de l'espace des $(a,b,c)$ se coupent en un seul point ; ou bien $(2,3,6)=3(1,1,1)-2(1,2,3)+1(1,4,9)$.
(d) Pour $x_1,x_2,x_3$ distincts, l'élimination de $\begin{pmatrix}1&x_1&x_1^2\\1&x_2&x_2^2\\1&x_3&x_3^2\end{pmatrix}$ donne les pivots $1$, $x_2-x_1$ et $(x_3-x_1)(x_3-x_2)$, tous non nuls.
\end{corrige}
\end{exo}
\begin{exo}[D][Une famille en bande]
Soit $K_n$ la matrice $n\times n$ avec des $2$ sur la diagonale, des $-1$ juste au-dessus et juste au-dessous, et des $0$ ailleurs. (a) Calculer $LU$ pour $n=3$ et $n=4$ et conjecturer les pivots $d_k$ et les multiplicateurs. (b) Démontrer la conjecture par récurrence. (c) On pose $M_{ij}=\frac{\min(i,j)\,(n+1-\max(i,j))}{n+1}$. Démontrer que $K_nM=I_n$.
\begin{indication}
(c) Fixe une colonne $j$ et regarde $i\mapsto M_{ij}$ : elle est affine pour $i\le j$ et pour $i\ge j$, et nulle en $i=0$ et $i=n+1$ si l'on prolonge la formule.
\end{indication}
\begin{corrige}
(a) Pour $n=3$ : $\ell_{21}=-\frac12$, ligne 2 $\to(0,\frac32,-1)$ ; $\ell_{32}=-\frac23$, ligne 3 $\to(0,0,\frac43)$. Pivots $2,\frac32,\frac43$ ; pour $n=4$, on ajoute $\ell_{43}=-\frac34$ et $d_4=\frac54$. Conjecture : $d_k=\frac{k+1}k$ et $\ell_{k+1,k}=-\frac{k}{k+1}$, aucun autre multiplicateur.
(b) Seul le coefficient $-1$ sous chaque pivot est à éliminer : la colonne $k$ n'a pas d'autre coefficient non nul sous la diagonale, et l'élimination de la ligne $k+1$ ne modifie que son coefficient diagonal. Si $d_k=\frac{k+1}k$, alors $\ell_{k+1,k}=\frac{-1}{d_k}=-\frac k{k+1}$, et la ligne $k+1$ perd $\ell_{k+1,k}$ fois la ligne $k$, dont le coefficient en colonne $k+1$ vaut $-1$ : $d_{k+1}=2-\ell_{k+1,k}\cdot(-1)=2-\frac k{k+1}=\frac{k+2}{k+1}$. Avec $d_1=2$, la récurrence est établie. En particulier tous les pivots sont non nuls et leur produit vaut $\frac21\cdot\frac32\cdots\frac{n+1}n=n+1$.
(c) Fixons $j$ et posons $m_i=M_{ij}$ pour $0\le i\le n+1$ (la formule donne $m_0=m_{n+1}=0$). Le coefficient $(i,j)$ de $K_nM$ est $-m_{i-1}+2m_i-m_{i+1}$. Pour $i<j$, les trois indices sont $\le j$ et $m_k=\frac{k(n+1-j)}{n+1}$ est affine en $k$ : la combinaison vaut $0$. De même pour $i>j$. Pour $i=j$ : $\frac{-(j-1)(n+1-j)+2j(n+1-j)-j(n-j)}{n+1}=\frac{(j+1)(n+1-j)-j(n-j)}{n+1}=\frac{n+1}{n+1}=1$. Donc $K_nM=I_n$ et $M=K_n^{-1}$ (exercice sur l'inverse unique). L'inverse d'une matrice creuse est ici entièrement pleine.
\end{corrige}
\end{exo}

\begin{bilan}
\item Multiplier des matrices de quatre façons ; utiliser le produit par blocs.
\item Rédiger Cauchy, Schwarz et l'inégalité triangulaire, avec les cas d'égalité.
\item Factoriser $PA=LU$ par éliminations explicites ; résoudre par descente et remontée ; inverser.
\item Démontrer \og inversible $\iff$ solution unique \fg{} et \og inversible $\iff$ colonnes libres \fg.
\item Reconnaître les ensembles de matrices stables par produit et par inverse.
\end{bilan}
